\subsection{映射的直接系统}

命题6. 设I是一个有向集，设 \((\mathbf{E}_{\alpha}, f_{\beta \alpha})\) 是相对于I的集合的直接系统，令 \(\mathbf{E} = \lim \mathbf{E}_{\alpha}\) 为其直接极限，且对每个 \(\alpha \in \mathbf{I}\) 让

\[
f _ {\alpha}: \mathbf {E} _ {\alpha} \rightarrow \mathbf {E}
\]

为典范映射。对每个 \(\alpha \in \mathbf{I}\) ，让 \(u_{\alpha}\) 为从 \(\mathbf{E}_{\alpha}\) 到一个集合 \(\mathbf{F}\) 中的映射，使得

\[
u _ {\beta} \circ f _ {\beta x} = u _ {\alpha}\tag{23}
\]

每当 \(\alpha \leqslant \beta\) .

那么：

\begin{enumerate}
\def\labelenumi{(\alph{enumi})}
\tightlist
\item
  存在唯一的映射 \(u\) （从 \(\mathbf{E}\) 到 \(\mathbf{F}\) 中），使得
\end{enumerate}

\[
u _ {\alpha} = u \circ f _ {\alpha}\tag{24}
\]

\[
{\text{对所有}} \alpha \in \mathbf {I}.
\]

\begin{enumerate}
\def\labelenumi{(\alph{enumi})}
\setcounter{enumi}{1}
\item
  \(u\) 是满射当且仅当 \(\mathbf{F}\) 是诸集合 \(u_{\alpha}(\mathbf{E}_{\alpha})\) 的并集。
\item
  \(u\) 是单射当且仅当对每个 \(\alpha \in I\) 这些关系 \(x \in E_{\alpha}\) , \(y \in E_{\alpha}\) , \(u_{\alpha}(x) = u_{\alpha}(y)\) 意味着存在 \(\beta \geqslant \alpha\) 使得 \(f_{\beta \alpha}(x) = f_{\beta \alpha}(y)\) .
\item
  采用第5号的记号，设 \(v\) 为从 \(G\) 到 \(F\) 中的、与 \(u_{\alpha}\) 在 \(E_{\alpha}\) 上（对于每一个 \(\alpha \in I\) ）一致的映射（第二章，第4节，第7号，命题8）。该假设意味着 \(v\) 与等价关系 \(R\) 相容（第二章，§6，第5号）；因此存在唯一的从 \(E = G / R\) 到 \(F\) 中的映射 \(u\) ，使得 \(v = u \circ f\) (同上引文)。
\item
  由于 \(\mathbf{E}\) 是诸 \(f_{\alpha}(\mathbf{E}_{\alpha})\) 的并集，关系 \(\mathbf{F} = \bigcup_{\alpha \in \mathbf{I}} u_{\alpha}(\mathbf{E}_{\alpha})\) 显然是 \(u\) 为满射的必要且充分条件。
\item
  根据第5号的引理1，E中的任意两个元素总可以写成以下形式 \(f_{\alpha}(x)\) 并且 \(f_{\alpha}(y)\) ，其中 \(x \in \mathbf{E}_{\alpha}\) 并且 \(y \in \mathbf{E}_{\alpha}\) ，对于合适的 \(\alpha \in \mathbf{I}\) 由该引理还可推出，关系 \(f_{\alpha}(x) = f_{\alpha}(y)\) 等价于存在 \(\beta \geqslant \alpha\) 使得 \(f_{\beta \alpha}(x) = f_{\beta \alpha}(y)\) 。由于 \(u_{\alpha}(x) = u(f_{\alpha}(x))\) 并且 \(u_{\alpha}(y) = u(f_{\alpha}(y))\) ，这就完成了证明。
\end{enumerate}

如果映射 \(u\) 是双射的，有时会滥用语言地说， \(F\) 是族 \((\mathbf{E}_{\alpha})\) 的直接极限。

备注。假设每个映射 \(f_{\beta \alpha}\) 是单射的。那么每一个 \(f_{\alpha}\) 是单射的，根据关系R的定义。在这种情况下，我们通常将 \(E_{\alpha}\) 与 \(f_{\alpha}(E_{\alpha})\) 等同起来，并因此将E视为诸 \(E_{\alpha}\) 的并集。反过来，设 \((F_{\alpha})_{\alpha \in I}\) 是集合 F 的子集的一个递增族，并假设 F 是该族的并集。如果 \(j_{\beta \alpha}\) 表示从 \(F_{\alpha}\) 到 \(F_{\beta}\) 中的典范单射（对于 \(\alpha \leqslant \beta\) ），那么由命题 6 可知，我们可以将 F 与族 \(F_{\alpha}\) 关于映射族 \((j_{\beta \alpha})\) 的直接极限等同起来，并且将从 \(F_{\alpha}\) 到 \(\lim F_{\alpha}\) 中的典范映射与从 \(F_{\alpha}\) 到 F 中的典范嵌入等同起来（对于每个 \(\alpha \in I\) ）。

推论1. 设 \((\mathbf{E}_{\alpha}, f_{\beta \alpha})\) 和 \((\mathbf{F}_{\alpha}, g_{\beta \alpha})\) 是相对于同一指标集 I 的两个集合直接系统；设 \(\mathbf{E} = \varinjlim \mathbf{E}_{\alpha}\) , \(\mathbf{F} = \varinjlim \mathbf{F}_{\alpha}\) ，并且对于每个 \(\alpha \in \mathbf{I}\) 让 \(f_{\alpha}\) （相应地， \(g_{\alpha}\) ）为从 \(\mathbf{E}_{\alpha}\) （相应地， \(\mathbf{F}_{\alpha}\) ）到 \(\mathbf{E}\) （相应地， \(\mathbf{F}\) ）中的典范映射。对于每一个 \(\alpha \in \mathbf{I}\) 让 \(u_{\alpha}\) 为从 \(\mathbf{E}_{\alpha}\) 到 \(\mathbf{F}_{\alpha}\) 中的一个映射，使得每当 \(\alpha \leqslant \beta\) 时，图表

\[
\begin{array}{c} \mathrm{E} _ {\alpha} \xrightarrow {u _ {\alpha}} \mathrm{F} _ {\alpha} \\ f _ {\beta \alpha} \Big \downarrow \qquad \qquad \qquad \qquad \qquad \qquad \qquad \qquad \qquad \Big \downarrow g _ {\beta \alpha} \\ \mathrm{E} _ {\beta} \xrightarrow [ u _ {\beta} ]{} \mathrm{F} _ {\beta} \end{array}
\]

是可交换的。则存在唯一的映射 \(u: \mathbf{E} \to \mathbf{F}\) ，使得对于每个 \(\alpha \in I\) ，图表

\[
\begin{array}{c} \mathrm{E} _ {\alpha} \xrightarrow {u _ {\alpha}} \mathrm{F} _ {\alpha} \\ f _ {\alpha} \Big \downarrow \qquad \qquad \qquad \qquad \Big \downarrow g _ {\alpha} \\ \mathrm{E} \xrightarrow {} \mathrm{F} \\ u \end{array}
\]

是交换的。

设 \(v_{\alpha} = g_{\alpha} \circ u_{\alpha}\) 。如果 \(\alpha \leqslant \beta\) ，那么由(22)我们有

\[
v _ {\beta} \circ f _ {\beta \alpha} = g _ {\beta} \circ u _ {\beta} \circ f _ {\beta \alpha} = g _ {\beta} \circ g _ {\beta \alpha} \circ u _ {\alpha} = g _ {\alpha} \circ u _ {\alpha} = v _ {\alpha}.
\]

因此，我们可以将命题6应用于诸映射 \(v_{\alpha}\) ，由此得出映射 \(u: \mathbf{E} \to \mathbf{F}\) 的存在性和唯一性，使得

\[
u \circ f _ {\alpha} = v _ {\alpha} = g _ {\alpha} \circ u _ {\alpha}
\]

对于所有 \(\alpha \in I\) .

一族映射 \(u_{\alpha} \colon \mathbf{E}_{\alpha} \to \mathbf{F}_{\alpha}\) ，若满足推论1的条件，则称为从 \((\mathbf{E}_{\alpha}, f_{\beta \alpha})\) 到 \((\mathbf{F}_{\alpha}, g_{\beta \alpha})\) 中的映射直接系统。推论1中定义的映射称为族 \((u_{\alpha})\) 的直接极限，并记作 \(u = \lim_{\rightarrow} u_{\alpha}\) 当没有歧义风险时。

推论2。设 \((\mathbf{E}_{\alpha}, f_{\beta \alpha})\) , \((\mathbf{F}_{\alpha}, g_{\beta \alpha})\) , \((\mathbf{G}_{\alpha}, h_{\beta \alpha})\) 是相对于I的三个集合直接系统。设 \(\mathbf{E} = \varinjlim \mathbf{E}_{\alpha}\) , \(\mathbf{F} = \varinjlim \mathbf{F}_{\alpha}\) , \(\mathbf{G} = \varinjlim \mathbf{G}_{\alpha}\) ，并且让 \(f_{\alpha}\) （相应地， \(g_{\alpha}, h_{\alpha}\) ）为从 \(\overrightarrow{\mathbf{E}_{\alpha}}\) （相应地， \(\mathbf{F}_{\alpha}, \overrightarrow{\mathbf{G}_{\alpha}}\) ）到 E（相应地，F，G）中的典范映射。如果 \((u_{\alpha})\) 和 \((v_{\alpha})\) 是两个映射直接系统 \(u_{\alpha}: \mathbf{E}_{\alpha} \to \mathbf{F}_{\alpha}\) , \(v_{\alpha}: \mathbf{F}_{\alpha} \to \mathbf{G}_{\alpha}\) ，那么映射 \(v_{\alpha} \circ u_{\alpha}: \mathbf{E}_{\alpha} \to \mathbf{G}_{\alpha}\) 构成一个映射直接系统，并且我们有

\[
\lim _ {\rightarrow} \left(v _ {\alpha} \circ u _ {\alpha}\right) = (\lim _ {\rightarrow} v _ {\alpha}) \circ (\lim _ {\rightarrow} u _ {\alpha}).\tag{25}
\]

因为如果我们令 \(w_{\alpha} = v_{\alpha} \circ u_{\alpha}\) ，那么对于 \(\alpha \leqslant \beta\) 我们拥有

\[
h _ {\beta \alpha} \circ w _ {\alpha} = \left(h _ {\beta \alpha} \circ v _ {\alpha}\right) \circ u _ {\alpha} = \left(v _ {\beta} \circ g _ {\beta \alpha}\right) \circ u _ {\alpha} = v _ {\beta} \circ \left(u _ {\beta} \circ f _ {\beta x}\right) = w _ {\beta} \circ f _ {\beta x},
\]

这表明 \((w_{\alpha})\) 是一个映射的直接系统。此外，如果 \(u = \varinjlim u_{\alpha}\) 并且 \(v = \varinjlim v_{\alpha}\) ，那么对所有 \(\alpha \in I\) 我们拥有

\[
(v \circ u) \circ f _ {\alpha} = v \circ (g _ {\alpha} \circ u _ {\alpha}) = h _ {\alpha} \circ (v _ {\alpha} \circ u _ {\alpha}),
\]

并且根据直接极限的唯一性，我们有 \(v \circ u = \lim_{\rightarrow} w_{\alpha}\) .

命题7。设 \((\mathbf{E}_{\alpha}, f_{\beta \alpha})\) 和 \((\mathbf{E}_{\alpha}', f_{\beta \alpha}')\) 为相对于I的两个集合直接系统，且对每个 \(\alpha \in I\) 让 \(u_{\alpha}\) 为从 \(\mathbf{E}_{\alpha}\) 到 \(\mathbf{E}_{\alpha}'\) 中的一个映射，使得各 \(u_{\alpha}\) 构成一个映射直接系统。设 \(u = \lim u_{\alpha}\) 。如果每个 \(u_{\alpha}\) 是单射（相应地，满射）则 \(u\) 是单射（相应地，满射）。

让 \(\mathbf{E} = \lim_{\rightarrow}\mathbf{E}_{\alpha},\mathbf{E}' = \lim_{\rightarrow}\mathbf{E}_{\alpha}'\) 并令 \(f_{\alpha}\colon \mathbf{E}_{\alpha}\to \mathbf{E},f_{\alpha}^{\prime}\colon \mathbf{E}_{\alpha}^{\prime}\to \mathbf{E}^{\prime}\) 为典范映射。假设每个 \(u_{\alpha}\) 是单射。要证明 \(u\) 是单射，根据命题 6，只需验证如果 \(x\in \mathbf{E}_{\alpha}\) 并且 \(y\in \mathbf{E}_{\alpha}\) 满足 \(f_{\alpha}^{\prime}(u_{\alpha}(x)) = f_{\alpha}^{\prime}(u_{\alpha}(y))\) ，那么存在 \(\beta \geqslant \alpha\) 使得 \(f_{\beta \alpha}(x) = f_{\beta \alpha}(y)\) 。现在该假设蕴含（第 6 号，引理 1）存在某个 \(\beta \geqslant \alpha\) 使得

\[
f _ {\beta \alpha} ^ {\prime} (u _ {\alpha} (x)) = f _ {\beta \alpha} ^ {\prime} (u _ {\alpha} (y)), \quad \text { i.e., } \quad u _ {\beta} (f _ {\beta \alpha} (x)) = u _ {\beta} (f _ {\beta \alpha} (y)),
\]

因此， \(f_{\beta \alpha}(x) = f_{\beta \alpha}(y)\) 因为 \(u_{\beta}\) 是单射的。

现在假设诸 \(u_{\alpha}\) 是满射的。于是我们有

\[
\mathrm{E} ^ {\prime} = \bigcup_ {\alpha} f _ {\alpha} ^ {\prime} (\mathrm{E} _ {\alpha} ^ {\prime}) = \bigcup_ {\alpha} f _ {\alpha} ^ {\prime} (u _ {\alpha} (\mathrm{E} _ {\alpha})) = \bigcup_ {\alpha} u (f _ {\alpha} (\mathrm{E} _ {\alpha})) = u \left(\bigcup_ {\alpha} f _ {\alpha} (\mathrm{E} _ {\alpha})\right) = u (\mathrm{E}).
\]

使用命题7的记号，设 \(\mathbf{M}_{\alpha}\) 为 \(\mathbf{E}_{\alpha}\) 的一个子集（对于每一个 \(\alpha \in \mathbf{I}\) ）；如果我们有 \(f_{\beta \alpha}(\mathbf{M}_{\alpha}) \in \mathbf{M}_{\beta}\) 每当 \(\alpha \leqslant \beta\) ，族 \((\mathbf{M}_{\alpha})_{\alpha \in \mathbf{I}}\) 被称为诸 \(\mathbf{E}_{\alpha}\) 的子集的直接系统。设 \(g_{\beta \alpha}\) （其中 \(\alpha \leqslant \beta\) ) 为从 \(\mathbf{M}_{\alpha}\) 到 \(\mathbf{M}_{\beta}\) 中的映射，其图形与 \(f_{\beta \alpha}\) 在 \(\mathbf{M}_{\alpha}\) 上的限制的图形相同。那么显然 \((\mathbf{M}_{\alpha}, g_{\beta \alpha})\) 是一个集合直接系统；并且命题7，应用于典范单射 \(j_{\alpha}: \mathbf{M}_{\alpha} \to \mathbf{E}_{\alpha}\) ，使我们能够将 \(\mathbf{M} = \lim \mathbf{M}_{\alpha}\) 与 \(\mathbf{E}\) 的一个子集通过单射 \(j = \lim j_{\alpha}\) 等同起来。

推论。设 \((\mathbf{E}_{\alpha}, f_{\beta\alpha})\) 和 \((\mathbf{E}_{\alpha}^{\prime}, f_{\beta\alpha}^{\prime})\) 是两个集合直接系统，设 \((u_{\alpha})\) 是一个映射直接系统 \(u_{\alpha}: E_{\alpha} \to E_{\alpha}^{\prime}\) ，并且让 \(u = \lim u_{\alpha}\) 。

\begin{enumerate}
\def\labelenumi{(\roman{enumi})}
\tightlist
\item
  让 \((\mathbf{M}_{\alpha})\) 是诸 \(\mathbf{E}_{\alpha}\) 的子集的一个直接系统。那么 \((u_{\alpha}(\mathbf{M}_{\alpha}))\) 是诸 \(\mathbf{E}_{\alpha}'\) 的子集的一个直接系统，并且我们有
\end{enumerate}

\[
\lim _ {\rightarrow} u _ {\alpha} (\mathbf {M} _ {\alpha}) = u (\lim _ {\rightarrow} \mathbf {M} _ {\alpha}).\tag{26}
\]

\begin{enumerate}
\def\labelenumi{(\roman{enumi})}
\setcounter{enumi}{1}
\tightlist
\item
  让 \((a_{\alpha}^{\prime})_{\alpha \in \mathbf{I}}\) 是一个族，使得 \(a_{\alpha}^{\prime} \in \mathrm{E}_{\alpha}^{\prime}\) 对于每一个 \(\alpha \in \mathbf{I}\) 并且 \(f_{\beta \alpha}^{\prime}(a_{\alpha}^{\prime}) = a_{\beta}^{\prime}\) 每当 \(\alpha \leqslant \beta\) 。则集合 \(-u_{\alpha}^{1}(a_{\alpha}^{\prime})\) 构成诸 \(\mathbf{E}_{\alpha}\) 的子集的一个直接系统，并且我们有
\end{enumerate}

\[
\lim _ {\rightarrow} \overline {{{u}}} _ {\alpha} ^ {1} (a _ {\alpha} ^ {\prime}) = \overline {{{u}}} ^ {1} (a ^ {\prime}),\tag{27}
\]

其中 \(a'\) 是 \(\lim_{\rightarrow} \mathbf{E}_{\alpha}'\) 中唯一的元素，它都是 \(a_{\alpha}'\) 的典范像（对每个 \(\alpha \in \mathbf{I}\) ）。

\begin{enumerate}
\def\labelenumi{(\roman{enumi})}
\item
  显然，各 \(u_{\alpha}(\mathbf{M}_{\alpha})\) 构成各集合 \(\mathbf{E}_{\alpha}'\) 的一个子集直接系统，并且可写成 \(u_{\alpha}(\mathbf{M}_{\alpha}) = v_{\alpha}(\mathbf{M}_{\alpha})\) ，其中 \(v_{\alpha}\) 是从 \(\mathbf{M}_{\alpha}\) 到 \(u_{\alpha}(\mathbf{M}_{\alpha})\) 上的映射，其图与 \(u_{\alpha}\) 在 \(\mathbf{M}_{\alpha}\) 上的限制之图相同。由于各 \(v_{\alpha}\) 都是满射，公式(26)便由命题7推出。
\item
  让 \(\mathbf{N}_{\alpha} = \overline{u}_{\alpha}^{1}(a_{\alpha}^{\prime})\) 。如果 \(\alpha \leqslant \beta\) 并且 \(x_{\alpha} \in \mathbf{N}_{\alpha}\) ，那么
\end{enumerate}

\[
u _ {\beta} (f _ {\beta \alpha} (x _ {\alpha})) = f _ {\beta \alpha} ^ {\prime} (u _ {\alpha} (x _ {\alpha}) = f _ {\beta \alpha} ^ {\prime} (a _ {\alpha} ^ {\prime}) = a _ {\beta} ^ {\prime};
\]

因此 \(f_{\beta \alpha}(x_{\alpha}) \in \mathbb{N}_{\beta}\) ，并且 \(\mathbb{N}_{\alpha}\) 因此构成 \(\mathbf{E}_{\alpha}\) 的子集的一个直接系统。采用命题7证明中的记号，考虑一个元素 \(x \in \lim \mathbb{N}_{\alpha}\) 。存在 \(\alpha \in \mathbb{I}\) 并且 \(x_{\alpha} \in \mathbb{N}_{\alpha}\) 使得 \(x = f_{\alpha}(x_{\alpha})\) ，因此 \(u(x) = u(f_{\alpha}(x_{\alpha})) = f_{\alpha}'(u_{\alpha}(x_{\alpha})) = f_{\alpha}'(a_{\alpha}') = a'\) 。反之，若 \(x \in \overline{u}^1(a')\) 并且如果 \(x = f_\alpha(x_\alpha)\) 对于某个 \(\alpha \in I\) 以及某个 \(x_\alpha \in E_\alpha\) ，那么我们有 \(a' = u(f_\alpha(x_\alpha)) = f_\alpha'(u_\alpha(x_\alpha)) = f_\alpha'(a_\alpha')\) 。因此（第5号，引理1）存在 \(\beta \geqslant \alpha\) 使得 \(f_{\beta\alpha}'(u_\alpha'(x_\alpha)) = f_{\beta\alpha}'(a_\alpha') = a_\beta'\) ；即， \(u_\alpha(f_{\beta\alpha}(x_\alpha)) = a_\beta'\) ，从而 \(f_{\beta\alpha}(x_\alpha) \in N_\beta\) 。由于 \(x = f_\beta(f_{\beta\alpha}(x_\alpha))\) ，由此可得 \(x \in \lim_{\rightarrow} N_\alpha\) .

备注。假设对每个 \(\alpha \in I\) , \(u_{\alpha}: E_{\alpha} \to E'\) 是一个映射，使得族 \((u_{\alpha})\) 满足 (23)。考虑相对于 I 的直接系统 \((E_{\alpha}, i_{\beta \alpha})\) ，其中 \(E_{\alpha}' = E'\) 对于所有 \(\alpha \in I\) ，以及 \(i_{\beta \alpha}\) 是 \(E'\) 的恒等映射。然后（第5号，例2） \(E'\) 可以典范地等同于 \(\lim E_{\alpha}'\) 。如果 \(u_{\alpha}\) 被视为从 \(E_{\alpha}\) 到 \(E_{\alpha}'\) 中的一个映射，那么 \((u_{\alpha})\) 是一个映射直接系统，并且由(24)定义的映射 \(u: E \to E'\) 被等同于该映射系统的直接极限。因此，通过滥用语言，我们写作 \(u = \lim u_{\alpha}\) 。

若J是I的一个子集，且J是有向的（关于诱导的预序），显然，由子族 \((\mathrm{E}_{\alpha})_{\alpha\in\mathbb{J}}\) 以及该族 \((f_{\beta\alpha})\) （其中 \(\alpha\leqslant\beta\) 并且 \(\alpha\in J\) 并且 \(\beta\in J\) ）构成的对是相对于J的集合直接系统；称其是通过将指标集限制为J而得到的。设E， \(E'\) 分别表示族 \((\mathrm{E}_{\alpha})_{\alpha\in\mathbb{I}}\) 和 \((\mathrm{E}_{\alpha})_{\alpha\in\mathbb{J}}\) 的直接极限，并且对于每个 \(\alpha\in I\) 让 \(f_{\alpha}:E_{\alpha}\to E\) 表示典范映射。则 \((f_{\alpha})_{\alpha\in\mathbb{J}}\) 是一个映射直接系统，因此 \(g=\lim f_{\alpha}\) 是从 \(E'\) 到 E 中的一个映射，称为典范的。此外，若 \(J'\) 是 J 的有向子集，若 \(E''\) 是族 \((\mathrm{E}_{\alpha})_{\alpha\in\mathbb{J}'}\) 的直接极限，并且如果

\[
g ^ {\prime}: \quad \mathbf {E} ^ {\prime \prime} \rightarrow \mathbf {E} ^ {\prime} \quad \text { 且 } \quad g ^ {\prime \prime}: \mathbf {E} ^ {\prime \prime} \rightarrow \mathbf {E}
\]

是典范映射，则由命题 6 立即得出

\[
g ^ {\prime \prime} = g \circ g ^ {\prime}.\tag{28}
\]

命题 8. 设 I 是有向集，设 \((\mathbf{E}_{\alpha}, f_{\beta \alpha})\) 是相对于I的一个集合直接系统，并令 \(\mathbf{E} = \lim \mathbf{E}_{\alpha}\) 为其直接极限。令 J 为 I 的共尾子集，并令 \(\mathbf{E}'\) 为由 \((\mathbf{E}_{\alpha}, f_{\beta \alpha})\) 通过将指标集限制为J而得到的集合直接系统的直接极限。则典范映射 \(g\) （从 \(\mathbf{E}'\) 到 E 中）是双射。

J 必然是一个有向集（§1，第10号）。我们将使用命题6的判据来证明g是双射的。单射的条件直接由定义和第5号引理1得出。为了证明g是满射的，我们注意到对于每个 \(\alpha \in J\) 我们拥有

\[
g (\mathbf {E} _ {\alpha}) = f _ {\alpha} (\mathbf {E} _ {\alpha}).
\]

现在，对每个 \(\beta \in I\) ，都存在 \(\gamma \in J\) 使得 \(\beta \leqslant \gamma\) ，由此 \(g(E_{\gamma}) \supset g(f_{\gamma \beta}(E_{\beta})) = f_{\beta}(E_{\beta})\) 。因此，E 是各集合 \(g(E_{\alpha})\) （当 \(\alpha\) 遍历 J 时）的并。

